\part{Packaging}
\chapter{Anatomy of the Package}
One Makefile target, \cmd{make install-pkg}, defines what ships. Both the \texttt{.deb} and the Arch
package are thin wrappers around it, so the two distributions cannot drift apart.

\section{From source to artifact}
\begin{figure}[ht]
\centering
\begin{tikzpicture}[x=1cm,y=1cm,scale=0.86,transform shape]
  \node[box,minimum width=2.6cm] (src) at (0,0) {source tree\\[-1pt]\scriptsize \texttt{kb-core kb-control-plane}\\[-3pt]\scriptsize \texttt{kb-op kb-checker kb-aads}};
  \node[box,minimum width=2.6cm] (build) at (4.2,0) {\texttt{make build}\\[-1pt]\scriptsize 5 stages};
  \node[box,minimum width=2.4cm] (bin) at (8,0) {\texttt{bin/} +\\[-1pt]\scriptsize \texttt{kbd\_sensor}};
  \node[box,minimum width=2.9cm] (stage) at (11.9,0) {\texttt{make install-pkg}\\[-1pt]\scriptsize \texttt{PREFIX=/usr DESTDIR=stage}};
  \node[dark,minimum width=3.8cm] (deb) at (11.9,-2.8) {\texttt{build-deb.sh}\\[-1pt]\scriptsize dpkg-deb $\to$ \texttt{.deb}};
  \node[dark,minimum width=3.8cm] (arch) at (6.6,-2.8) {\texttt{makepkg}\\[-1pt]\scriptsize PKGBUILD $\to$ \texttt{.pkg.tar.zst}};
  \draw[flow] (src) -- (build); \draw[flow] (build) -- (bin); \draw[flow] (bin) -- (stage);
  \draw[flow] (stage) -- (deb);
  \draw[flow] (stage.south west) -- ++(0,-0.9) -| (arch.north);
  \node[lbl] at (3.2,-1.6) {build on the TARGET distro};
  \draw[decorate,decoration={brace,mirror,amplitude=5pt},line width=0.8pt] (-1.5,-0.9) -- (9.2,-0.9);
\end{tikzpicture}
\caption{The packaging pipeline. The brace marks the part that must run on the distribution you are
targeting, because the binaries link the build host's C library.}
\end{figure}

\begin{kbwarn}[BUILD ON THE DISTRIBUTION YOU SHIP TO]
The Go, Rust and C binaries are dynamically linked against the \emph{build host's} glibc. The Arch host that
builds the Arch package has glibc 2.44; Ubuntu 24.04 has 2.39. A binary built on the former will not start
on the latter (``version \texttt{GLIBC\_2.4x} not found''). That is why the \texttt{.deb} is built on the
Ubuntu machine and the Arch package on Arch. Building on the \emph{older} distribution always works
forward; the reverse does not.
\end{kbwarn}

\section{The five build stages}
\cmd{make build} runs these in order and puts every product into the top-level \file{bin/}
(the sensor stays in \file{kb-core/build/}).
\begin{center}\small
\begin{tabularx}{\textwidth}{@{}L{1.2cm}L{3.2cm}Y@{}}
\toprule
\textbf{Stage} & \textbf{Output} & \textbf{What happens}\\
\midrule
1/5 & \file{kb-core/build/kbd_sensor} & \cmd{clang} compiles the BPF programs and userspace loader, \cmd{bpftool} generates the CO-RE skeletons, the vendored \file{libbpf/src/libbpf.a} is linked statically. Hardened with \texttt{-D\_FORTIFY\_SOURCE=3}, stack protector, \texttt{-fPIE}, full RELRO.\\
2/5 & \file{bin/kbd} & \texttt{go build -a} of \file{kb-control-plane/cmd/kbd}.\\
3/5 & \file{bin/kb-checker} & \texttt{cargo build --release} of the watchdog (built even though it is kept down).\\
4/5 & \file{bin/kbctl} & \texttt{go build -a} of the CLI.\\
5/5 & \file{bin/mcp-server}, \file{bin/kbopd} & Builds the MCP server; runs \texttt{npm ci} (only if \file{node_modules} is absent) and \texttt{npm run build} for the dashboard; then \texttt{cargo build --release} of \file{kb-tui}, which \emph{embeds} the dashboard's \file{dist/} at compile time and is copied to \file{bin/kbopd}.\\
\bottomrule
\end{tabularx}
\end{center}
\begin{kbnote}[ORDER MATTERS IN STAGE 5]
The dashboard must be built \emph{before} \cmd{cargo} compiles \file{kb-tui}. The Rust code embeds
\file{kb-dashboard/dist/} with \texttt{rust-embed}; if the folder is missing the compile fails with
``folder ... does not exist''. The Makefile already enforces the order.
\end{kbnote}

\section{The Makefile contract}
\begin{center}\small
\begin{tabularx}{\textwidth}{@{}L{3.3cm}L{3.9cm}Y@{}}
\toprule
\textbf{Variable} & \textbf{Default} & \textbf{Meaning}\\
\midrule
\texttt{PREFIX} & \texttt{/usr/local} & Packages pass \texttt{/usr}.\\
\texttt{BINDIR} & \texttt{\$(PREFIX)/bin} & Where binaries land; also substituted into unit files and the SSH wrapper.\\
\texttt{LIBDIR} & \texttt{\$(PREFIX)/lib} & \file{kb/provision.sh} and \file{kb/aads/}.\\
\texttt{SYSCONFDIR} & \texttt{/etc} & Config root (\file{/etc/kb}).\\
\texttt{UNITDIR} & \texttt{/etc/systemd/system} for \texttt{/usr/local}, else \file{$(PREFIX)/lib/systemd/system} & Where unit files go.\\
\texttt{DESTDIR} & empty & Staging root, prepended to every install path.\\
\texttt{GOARCH} & \texttt{go env GOARCH} & Target architecture of the Go binaries.\\
\bottomrule
\end{tabularx}
\end{center}
\begin{center}\small
\begin{tabularx}{\textwidth}{@{}L{3.3cm}Y@{}}
\toprule
\textbf{Target} & \textbf{Effect}\\
\midrule
\texttt{make build} & The five stages above.\\
\texttt{make install} & Installs the six binaries to \texttt{\$(DESTDIR)\$(BINDIR)}. Historical behaviour: \texttt{sudo make install} still goes to \file{/usr/local/bin}.\\
\texttt{make install-pkg} & \texttt{install} plus the wrapper, \cmd{kbagents}, units (path-rewritten), operator \cmd{sshd} config, default \file{/etc/kb/*.yaml}, \file{provision.sh}, the \cmd{kb-aads} sources, and docs.\\
\texttt{make clean} & Removes build products.\\
\bottomrule
\end{tabularx}
\end{center}

\begin{kbwarn}[THE PACKAGE USES EXPLICIT FILE LISTS ON PURPOSE]
The source tree contains development-only \file{ssh_host_ed25519_key} and \file{authorized_keys} files under
\file{kb-core/} and \file{kb-control-plane/}. They are git-ignored, but they exist on disk. \cmd{install-pkg} names every
file it ships and never copies a directory wholesale (except the \cmd{kb-aads} Python packages, which are
listed folder by folder). If you extend the package, keep it that way and re-run the secrets check in
Section~\ref{sec:secrets}.
\end{kbwarn}

\section{What the unit files are rewritten to}
The sources in \file{packaging/systemd/} use the development path \file{/usr/local/bin}. During
\cmd{install-pkg}, \cmd{sed} rewrites it to \texttt{\$(BINDIR)} (\file{/usr/bin}) in every unit and in the SSH wrapper, and
rewrites the operator \cmd{sshd} config path from \file{/etc/ssh/sshd_config.d/kb-operator.conf} to
\file{/etc/kb/sshd_kb-operator.conf}.

\begin{kbnote}[WHY THE SSHD CONFIG MOVED]
Ubuntu's main \cmd{sshd} reads \texttt{Include /etc/ssh/sshd\_config.d/*.conf}. A file there containing
\texttt{Port 2222} is therefore also loaded by the \emph{host's own} \cmd{sshd}, which grabs port 2222 first;
the dedicated \texttt{sshd@kb-operator} instance then dies with \texttt{Address already in use} and restarts
forever. This was observed on the live Ubuntu VM. The packages avoid it by keeping the file under
\file{/etc/kb/}, where only the dedicated unit reads it.
\end{kbnote}

\section{Package metadata and maintainer scripts}
\begin{center}\small
\begin{tabularx}{\textwidth}{@{}L{3.1cm}L{3.2cm}Y@{}}
\toprule
\textbf{Item} & \textbf{Ubuntu \texttt{.deb}} & \textbf{Arch}\\
\midrule
Source of truth & \file{packaging/deb/} & \file{packaging/arch/PKGBUILD}\\
Runtime deps & \texttt{libc6, systemd, openssh-server, libelf1t64, zlib1g} & \texttt{glibc systemd openssh libelf zlib}\\
Optional & \texttt{python3, python3-venv} (Recommends) & \texttt{python} (optdepends)\\
Config preserved on upgrade & \file{conffiles} (five files) & \texttt{backup=(...)} array (five files)\\
After install & \texttt{postinst} runs \file{provision.sh}, \cmd{systemctl daemon-reload} & \file{.install} \texttt{post\_install} does the same and prints next steps\\
Before removal & \texttt{prerm} disables and stops the units & \texttt{pre\_remove} does the same\\
After removal & \texttt{postrm} reloads systemd; keeps \file{/etc/kb}, state, users & \texttt{post\_remove} reloads systemd; same retention\\
Build flags & plain \texttt{make build} & \texttt{!lto}, neutral \texttt{CFLAGS/LDFLAGS/RUSTFLAGS}\\
\bottomrule
\end{tabularx}
\end{center}
\textbf{Nothing is enabled or started on install.} \cmd{kb-checker} can legitimately run a network lockdown;
a package must never arm that by surprise. You start the stack yourself with \cmd{kbctl up}
(Chapter~\ref{ch:start}).

\section{What is deliberately not in the package}
\begin{itemize}
\item \textbf{Development keys and datasets.} No private key, no \file{authorized_keys} from the tree, no \file{scripts/dataset}, no zip bundles.
\item \textbf{The Python virtualenv.} The \cmd{kb-aads} environment (Ray, PyTorch) is multiple gigabytes and CPU-specific; you create it on the target (Chapter~\ref{ch:aads}).
\item \textbf{Build outputs of other subsystems} such as \file{target/} directories and \file{node_modules}.
\end{itemize}
The \cmd{kb-aads} source tree \emph{does} include one trained policy checkpoint, \file{marl/checkpoints/containment_ppo/}
(16 \file{.pkl} files, 1.8~MB). It is shipped because the containment agent loads it at runtime
(\texttt{RLModule.from\_checkpoint} in \file{agents/containment.py}); without it the agent reports
\texttt{model\_loaded: False} and recommends no containment. Pickle files execute code when loaded, so:
\begin{itemize}
\item They ship root-owned and not world-writable, under \file{/usr/lib/kb/aads/}, which only root may modify.
\item The package carries an integrity manifest. Verify it any time, and before every start if you are cautious:
\begin{lstlisting}[style=console]
$ cd /usr/lib/kb/aads/marl/checkpoints && sha256sum -c SHA256SUMS
containment_ppo/learner_group/state.pkl: OK
...
\end{lstlisting}
\item The unused \texttt{jury\_ppo} checkpoint and the development \file{session.log} are \emph{not} shipped.
\end{itemize}
The manifest protects against accidental damage and casual tampering; an attacker with root can replace both
files, which is why \cmd{kb-checker}'s eventual file-integrity role matters (Chapter~\ref{ch:checker}).
